% ============================================================
% 《算法的低语·数学卷》独立分册 preamble
% 基于主册 preamble.tex 精简（去除主册封面与变更日志，新增数学卷封面）
% ============================================================

\usepackage{fontspec}
\usepackage[fontset=fandol]{ctex}
\usepackage{geometry}
\geometry{
  a4paper,
  top=2.7cm, bottom=2.7cm, left=2.8cm, right=2.8cm,
  headheight=1.2cm, footskip=1.4cm
}
\usepackage{amsmath, amssymb, mathtools}
\usepackage{graphicx}
\usepackage{float}
\usepackage{xcolor}
\usepackage{hyperref}
\usepackage[open,openlevel=1]{bookmark}
\usepackage{eso-pic}
% 恢复 secnumdepth，让 hyperref 能给每个 section/subsection 独立 anchor（不显示编号）
\AtBeginDocument{\setcounter{secnumdepth}{3}}
% 但隐藏实际编号显示
\renewcommand{\thesection}{}
\renewcommand{\thesubsection}{}
\renewcommand{\thesubsubsection}{}
\usepackage{tikz}
\usepackage{fancyhdr}
\usepackage{titlesec}
\usepackage{enumitem}
\usepackage{booktabs}
\usepackage{array}
\usepackage{calc}
\usepackage[most]{tcolorbox}

% ---- 颜色（与主册一致）----
\definecolor{teal1}{RGB}{0,77,89}
\definecolor{teal2}{RGB}{15,118,135}
\definecolor{teal3}{RGB}{38,162,180}

% ---- 占位符：用三角形作为导读盒标题前缀 ----
\newcommand{\faIconStub}{\textbf{\textcolor{teal2}{▶}}}

% ---- 核心概念横条：每节标题下方,3-6 个关键名词,紧贴节标题不抢戏 ----
\newtcolorbox{keytermsbox}{
  enhanced,
  colback=teal3!4,
  colframe=teal3!25,
  boxrule=0.4pt,
  arc=2pt,
  left=8pt, right=8pt, top=4pt, bottom=4pt,
  before skip=2pt, after skip=8pt,
  fontupper=\small
}
\newcommand{\keyterms}[1]{%
  \begin{keytermsbox}%
    \textbf{\color{teal1}核心概念}\hspace{0.6em}\textcolor{teal3!50}{|}\hspace{0.6em}%
    \color{teal1}#1%
  \end{keytermsbox}%
}

% ---- 小节导读盒：「这一节在讲什么」----
\newtcolorbox{sectionintuition}{
  enhanced,
  colback=teal3!6,
  colframe=teal2,
  boxrule=0.6pt,
  arc=3pt,
  left=10pt, right=10pt, top=8pt, bottom=8pt,
  before skip=8pt, after skip=10pt,
  fonttitle=\bfseries\color{teal1},
  title={\faIconStub\,这一节在讲什么},
  attach title to upper,
  after title={\par\vspace{4pt}}
}

% ---- 超链接 ----
\hypersetup{
  colorlinks=true,
  linkcolor=teal2,
  urlcolor=teal2,
  citecolor=teal2,
  pdftitle={算法的低语·数学卷 v1.7},
  pdfauthor={大队长}
}

% ---- 页眉页脚（与主册风格统一）----
\pagestyle{fancy}
\fancyhf{}
\fancyhead[L]{\small\color{gray}\itshape 算法的低语·数学卷}
\fancyhead[R]{\small\color{gray}\itshape 从公式到推导}
\fancyfoot[C]{\small\color{gray}\thepage}
\fancyfoot[R]{\small\color{gray}\itshape 大队长出品}
\renewcommand{\headrulewidth}{0.3pt}
\renewcommand{\headrule}{\color{teal3!40}\hrule width\headwidth height\headrulewidth}

% ---- 水印（每页对角线「大队长出品」，仅正文页；封面与版权页跳过）----
% 以页面中心为锤点旋转，避免超出页面边界
\AddToShipoutPictureBG{%
  \ifnum\value{page}>2
    \AtPageCenter{%
      \makebox(0,0){%
        \rotatebox{45}{\color{gray!7}\fontsize{72}{72}\selectfont\bfseries 大队长出品}%
      }%
    }%
  \fi
}

% ---- 标题样式 ----
\titleformat{\section}
  {\Large\bfseries\color{teal1}}
  {\thesection}{0.6em}{}
  [\vspace{0.2em}{\color{teal3!50}\hrule height 0.5pt}]

\titleformat{\subsection}
  {\large\bfseries\color{teal2}}
  {\thesubsection}{0.5em}{}

\titleformat{\subsubsection}
  {\normalsize\bfseries\color{teal2}}
  {\thesubsubsection}{0.4em}{}

% ---- 段落 ----
\setlength{\parindent}{0pt}
\setlength{\parskip}{0.8em}
\linespread{1.45}

% 列表间距
\setlist{itemsep=0.35em, parsep=0.25em, topsep=0.4em}

% 节/小节标题上下间距
\titlespacing*{\section}
  {0pt}{2.2em plus 0.4em minus 0.2em}{1.0em}
\titlespacing*{\subsection}
  {0pt}{1.6em plus 0.3em minus 0.2em}{0.7em}
\titlespacing*{\subsubsection}
  {0pt}{1.2em plus 0.3em minus 0.2em}{0.5em}

% ---- 数学卷封面（不用 titlepage，避免页码计数器重置导致书签偏移）----
\newcommand{\makemathcoverpage}{%
  \begingroup
    \thispagestyle{empty}
    \pagecolor{teal1}
    \begin{center}
      \vspace*{1.8cm}
      {\color{white!80}\rule{6cm}{0.6pt}\par}
      \vspace{0.8cm}
      {\fontsize{48}{56}\selectfont\bfseries\color{white} 算法的低语\par}
      \vspace{0.5cm}
      {\fontsize{28}{32}\selectfont\bfseries\color{teal3!40!white} 数\enspace 学\enspace 卷\par}
      \vspace{0.3cm}
      {\fontsize{15}{18}\selectfont\itshape\color{teal3!60!white} 从公式到推导\par}
      \vspace{0.8cm}
      {\color{white!80}\rule{6cm}{0.6pt}\par}

      \vspace{1.8cm}
      \begin{minipage}{0.72\textwidth}
        \itshape\color{white!75}\small\setlength{\parskip}{0.4em}
        本册是《算法的低语》主本的姊妹册，镜像主本7章结构，对每个AI概念背后的数学给出三层呈现：

        \begin{itemize}\setlength\itemsep{0.15em}\color{white!75}
          \item \textbf{直觉}（150–250字）：回顾概念由来与几何含义
          \item \textbf{严格}（400–800字 + 公式块）：定义、推导步骤、矩阵维度
          \item \textbf{全编号}：426条公式逐一编号，叙述以「式(X.Y.N)」明确引用
        \end{itemize}

        面向已具备线性代数、微积分、概率统计本科基础的进阶读者。建议与主本《算法的低语》配套阅读。
      \end{minipage}

      \vfill
      {\color{white}\Large\bfseries 大队长\enspace 出品\par}
      \vspace{0.25em}
      {\color{teal3!55!white}\small\href{mailto:fqsx@mail.ustc.edu.cn}{fqsx@mail.ustc.edu.cn}\par}
      \vspace{0.4em}
      {\color{teal3!45!white}\footnotesize v1.7\;·\;公式全编号版\;·\;52 节\,/\,426 条编号式\par}
      \vspace{0.4cm}
    \end{center}
    \clearpage
    \nopagecolor
  \endgroup
}
